\documentclass[12pt,a4paper]{report}
\usepackage[utf8]{inputenc}
\usepackage[brazil]{babel}
\usepackage{amsmath}
\usepackage{amsfonts}
\usepackage{amssymb}
\usepackage{enumitem}
\usepackage{fancyhdr}
\usepackage{geometry}
\usepackage{graphicx}
\usepackage{multicol}
\usepackage{tcolorbox}
\usepackage{parskip}
\usepackage{setspace}
\usepackage{tikz}
\usepackage{pgfplots}
\pgfplotsset{compat=1.15}
\usepackage{verbatim}
\usetikzlibrary{arrows.meta, calc}
\setstretch{1.2}
\tcbuselibrary{skins, breakable}
\geometry{top=2cm, bottom=2cm, left=2.5cm, right=2.5cm}
\usepackage[colorlinks=true, linkcolor=blue, urlcolor=blue, citecolor=blue]{hyperref}

\usepackage{chngcntr}
\counterwithout{footnote}{chapter}

\pagestyle{fancy}
\fancyhf{} 
\fancyhead[C]{\leftmark} 
\fancyfoot[C]{\thepage} 

% Estilo de caixa para destaque conceitual
\newtcolorbox{destaque}[1][]{
	enhanced,
	colback=gray!5,
	colframe=black!70,
	boxrule=1pt,
	arc=4pt,
	left=10pt,
	right=10pt,
	top=8pt,
	bottom=8pt,
	fonttitle=\bfseries,
	code={%
		\ifstrempty{#1}{%
			\pgfkeysalso{notitle}%
		}{%
			\pgfkeysalso{title={#1}}%
		}%
	}
}

\begin{document}
	
	\begin{titlepage}
		\centering
		
		\includegraphics[width=8cm]{ifes.png}\\[1cm]
		
		{\Large \textbf{Instituto Federal do Espírito Santo}}\\[0.5cm]
		{\large \textbf{Licenciatura em Matemática}}\\[0.5cm]
		{\large Introdução ao Estudo das Funções}
		
		\vspace{\fill}
		
		{\Huge \textbf{Função Exponencial}}\\[0.5cm]
		{\Large Material de Apoio}
		
		\vspace{\fill}
		
		{\Large Professor Alex Jordane}
		
		\vspace{\fill}
		
		{\large Vitória -- ES}\\[0.2cm]
		{\large 2026}
		
	\end{titlepage}
	
	\clearpage
%	\pagenumbering{roman} % Define a numeração em algarismos romanos minúsculos (i, ii, iii, ...)
	\setcounter{page}{1}  % Inicia a contagem pré-textual a partir de i
	
	\chapter*{Apresentação}
	\addcontentsline{toc}{chapter}{Apresentação}
	\thispagestyle{empty}
	
	Neste material, dedicamo-nos ao estudo aprofundado das funções exponenciais e suas aplicações analíticas. Iniciamos com as definições formais e as condições de existência da base, articuladas às propriedades estruturais da potenciação. Em seguida, investigamos a representação geométrica da curva e suas características centrais: o comportamento assintótico horizontal, a interseção com o eixo das ordenadas, o domínio, o conjunto imagem, a monotonicidade estrita (crescimento e decrescimento), o estudo de sinais e os efeitos provocados pelos parâmetros de translação e dilatação.
	
	Após a sistematização desses conceitos fundamentais, ampliamos o percurso para o tratamento das equações e inequações exponenciais. Detalhamos os métodos da redução a uma base comum e da mudança de variável, dedicando atenção especial ao exame de potências com bases variáveis. Para consolidar essas reflexões, disponibilizamos conjuntos de exercícios propostos seguidos de suas respectivas resoluções comentadas passo a passo, pensadas para apoiar os momentos de consulta, estudo autônomo e trabalho colaborativo em sala de aula.
	
	Bons estudos!
	
	\vspace{\fill}
	
	{\footnotesize
		\texttt{Esta é uma segunda versão do material, atualizada em 03/11/2026.}\\
		\texttt{Este material foi elaborado em colaboração com a inteligência artificial Gemini, \\ versão 3.1 Pro.} \\		
		\texttt{O compartilhamento e a reprodução deste material são permitidos, desde que a fonte \\ e a autoria sejam citadas.}  \\		
		\texttt{Esta obra está licenciada com uma Licença Atribuição-NãoComercial-SemDerivações 4.0.}
		\vspace{1cm}
		\begin{center}
			\includegraphics[width=4cm]{CC.png}
	\end{center}}
	
	\clearpage
	\chapter*{Convenções e Notações Matemáticas}
	\addcontentsline{toc}{chapter}{Convenções e Notações Matemáticas}
	\thispagestyle{empty}
	
	Para assegurar a clareza na leitura e no acompanhamento das deduções, destacamos as principais convenções e símbolos adotados ao longo deste material:
	
	\vspace{0.3cm}
	
	\begin{destaque}[Intervalos Reais (Notação de Bourbaki)]
		Para representar intervalos numéricos na reta real, adotamos a convenção de colchetes voltados para fora para indicar extremidades abertas:
		\begin{itemize}[leftmargin=1.5cm, label=--]
			\item $[a, b] = \{x \in \mathbb{R} \mid a \leqslant x \leqslant b\}$ \quad (Intervalo fechado)
			\item ${]a, b[} = \{x \in \mathbb{R} \mid a < x < b\}$ \quad (Intervalo aberto)
			\item ${[a, b[} = \{x \in \mathbb{R} \mid a \leqslant x < b\}$ \quad (Intervalo semiaberto à direita)
			\item ${]a, b]} = \{x \in \mathbb{R} \mid a < x \leqslant b\}$ \quad (Intervalo semiaberto à esquerda)
			\item ${[a, +\infty[} = \{x \in \mathbb{R} \mid x \geqslant a\}$ \quad e \quad ${]-\infty, b]} = \{x \in \mathbb{R} \mid x \leqslant b\}$
			\item ${]a, +\infty[} = \{x \in \mathbb{R} \mid x > a\}$ \quad e \quad $]-\infty, b[ = \{x \in \mathbb{R} \mid x < b\}$
		\end{itemize}
	\end{destaque}
	
	\vspace{0.4cm}
	
	\section*{Símbolos e Conjuntos Numéricos}
	
	\begin{itemize}[leftmargin=1.5cm, label=$\bullet$]
		\item $\mathbb{R}$: Conjunto dos números reais.
		\item $\mathbb{R}^*$: Conjunto dos números reais não nulos, isto é, $\mathbb{R} - \{0\}$.
		\item $\mathbb{R}_+$ e $\mathbb{R}^*_+$: Conjunto dos números reais não negativos e estritamente positivos (${]0, +\infty[}$), respectivamente.
		\item $\text{Dom}(f)$ e $\text{Im}(f)$: Domínio e conjunto imagem de uma função $f$.
		\item $S$: Conjunto solução de uma equação ou inequação.
		\item $\iff$: Equivalência lógica (``se, e somente se'').
		\item $\implies$: Implicação lógica direta.
	\end{itemize}
	
	\clearpage
	
	\tableofcontents
	\thispagestyle{empty}
	
	\clearpage
	%\pagenumbering{arabic} % Retorna à numeração arábica tradicional (1, 2, 3, ...)
	%\setcounter{page}{1}   % Reinicia a contagem para que o Capítulo 1 comece na página 1
		
	\chapter{Fundamentos Algébricos e Condições de Existência}
	\label{sec:algebricas_exp}
	
	Diferentemente das funções polinomiais, nas quais a variável independente situa-se na base e os expoentes são números naturais fixos, na função exponencial a variável assume a posição de expoente. Essa mudança estrutural exige uma investigação cuidadosa sobre o domínio e a natureza da base.
	
	\section{Definição e Condição de Existência}
	\label{subsec:cond_existencia}
	
	\begin{destaque}[Função Exponencial Elementar]
		\label{dest:def_exp}
		Chama-se \textbf{função exponencial elementar} a aplicação $f: \mathbb{R} \to \mathbb{R}$ dada pela lei:
		$$f(x) = a^x$$
		onde a base $a$ satisfaz necessariamente:
		$$a \in \mathbb{R}, \quad a > 0 \quad \text{e} \quad a \neq 1$$
	\end{destaque}
	
	A exigência de que $a > 0$ e $a \neq 1$ decorre de três impedimentos algébricos e funcionais em $\mathbb{R}$:
	
	\begin{enumerate}
		\item \textbf{Se $a < 0$:} A potência $a^x$ torna-se indefinida para infinitos valores reais no domínio. Por exemplo, se $a = -4$ e $x = \frac{1}{2}$, teríamos $(-4)^{1/2} = \sqrt{-4} \notin \mathbb{R}$. Para garantir que a função tenha domínio em toda a reta real $\mathbb{R}$ sem descontinuidades ou saltos imaginários, impõe-se $a > 0$.
		\item \textbf{Se $a = 0$:} A expressão $0^x$ gera indeterminação matemática para $x = 0$ ($0^0$) e é indefinida para valores negativos de $x$ (pois envolveria divisão por zero, como em $0^{-2} = \frac{1}{0^2}$). Para $x > 0$, resultaria em uma função identicamente nula, desprovida da dinâmica exponencial.
		\item \textbf{Se $a = 1$:} A relação $f(x) = 1^x = 1$ converte-se em uma função constante cuja reta é horizontal e paralela ao eixo das abscissas. Embora bem definida em $\mathbb{R}$, ela não compartilha as propriedades essenciais da família exponencial, tais como a estrita monotonicidade e a inversibilidade.
	\end{enumerate}
	
	\section{Propriedades Estruturais Operatórias}
	\label{subsec:propriedades}
	
	A extensão da potenciação para expoentes reais preserva as regras fundamentais da aritmética, fundamentando a manipulação algébrica das equações e funções exponenciais. Esse movimento começa com os números naturais e se estende até os números reais.

	\subsection{Extensão do Conceito de Potência}
	
	A potência $a^n$ é construída por meio de sucessivas extensões numéricas para o expoente:
	
	\begin{itemize}
		\item \textbf{Expoente Natural ($n \in \mathbb{N}^*$):} $a^1 = a$ \quad e \quad $a^n = a^{n-1} \cdot a \quad (n \geqslant 2, \, a \in \mathbb{R})$.
		\item \textbf{Expoente Nulo e Inteiro Negativo ($n \in \mathbb{Z}$):} $a^0 = 1$ \quad e \quad $a^{-n} = \dfrac{1}{a^n} \quad (a \in \mathbb{R}^*)$.
		\item \textbf{Expoente Racional ($q = \frac{m}{n} \in \mathbb{Q}$):} $a^{\frac{m}{n}} = \sqrt[n]{a^m} \quad (a \in \mathbb{R}^*_+, \, m \in \mathbb{Z}, \, n \in \mathbb{N}^*)$.
		\item \textbf{Expoente Real ($x \in \mathbb{R}$):} Definido por limites ou aproximações de expoentes racionais, exigindo $a \in \mathbb{R}^*_+$.
	\end{itemize}
	
	\subsection{Propriedades Operatórias Fundamentais}
	
	Com a base positiva garantida ($a, b \in \mathbb{R}^*_+$), todas as propriedades aritméticas preservam sua validade para quaisquer expoentes reais:
	
	\begin{destaque}[{Propriedades das Potências ($a, b \in \mathbb{R}^*_+$ e $m, n \in \mathbb{R}$)}]
		\begin{align*}
			a^m \cdot a^n &= a^{m+n} & \frac{a^m}{a^n} &= a^{m-n} \\[0.3cm]
			(a^m)^n &= a^{m \cdot n} & (a \cdot b)^n &= a^n \cdot b^n \\[0.3cm]
			\left(\frac{a}{b}\right)^n &= \frac{a^n}{b^n} & a^{-n} &= \frac{1}{a^n} = \left(\frac{1}{a}\right)^n
		\end{align*}
	\end{destaque}
	
	\section{A Forma Geral Parametrizada}
	\label{subsec:forma_geral_exp}
	
	Em aplicações empíricas e modelagens de fenômenos contínuos, a função exponencial aparece frequentemente acompanhada de fatores de escala, translação e contração temporal:
	
	\begin{destaque}[Forma Geral Parametrizada]
		\label{dest:forma_geral_param}
		$$f(x) = k \cdot a^{c(x - x_0)} + d$$
		com $k, c, x_0, d \in \mathbb{R}$, $k \neq 0$, $c \neq 0$, $a > 0$ e $a \neq 1$.
	\end{destaque}
	
	Cada parâmetro atua de forma direta sobre a geometria da curva:
	\begin{itemize}
		\item $k$: Fator de dilatação vertical e reflexão em relação ao eixo $x$ (quando $k < 0$).
		\item $c$: Fator de escala horizontal que comprime ($|c| > 1$) ou dilata ($0 < |c| < 1$) a curva.
		\item $x_0$: Translação horizontal: deslocamento para a direita se $x_0 > 0$, ou para a esquerda se $x_0 < 0$).
		\item $d$: Translação vertical: desloca a curva verticalmente para cima se $d>0$, ou para baixo se $d<0$.
	\end{itemize}
	
	\chapter{Representação Gráfica e Geometria da Curva}
	\label{sec:grafico_exp}
	
	A representação gráfica da função exponencial elementar $f(x) = a^x$ assume duas configurações, dependendo de a base ser maior do que um ou estar compreendida entre zero e um.
	
	\section{Comportamento da Base: $a > 1$ versus $0 < a < 1$}
	\label{subsec:base_comparacao}
	
	\begin{itemize}
		\item \textbf{Caso $a > 1$ (Crescimento Exponencial):} A função cresce indefinidamente à medida que $x$ avança no semi-eixo positivo, acelerando sua curvatura. À medida que $x \to -\infty$, as imagens aproximam-se cada vez mais de zero, sem nunca tocá-lo. Dizemos então que $y=0$ é uma \textbf{assíntota horizontal}.
		\item \textbf{Caso $0 < a < 1$ (Decaimento Exponencial):} A função decresce continuamente. À medida que $x \to +\infty$, a curva se aproxima do eixo horizontal, sem nunca tocá-lo (\textbf{assíntota horizontal}); quando $x \to -\infty$, os valores assumem magnitudes arbitrariamente grandes.
	\end{itemize}
	
	\begin{center}
		\begin{tikzpicture}[>=Stealth, scale=0.95]
			% Eixos
			\draw[->, line width=0.8pt] (-3.5,0) -- (3.5,0) node[right] {$x$};
			\draw[->, line width=0.8pt] (0,-0.6) -- (0,5.2) node[above] {$y$};
			\node[below left, font=\footnotesize] at (0,0) {$0$};
			
			% Curva a > 1: f(x) = 2^x
			\draw[domain=-3.2:2.3, smooth, line width=1.3pt, blue!85!black] plot ({\x}, {2^(\x)});
			\node[right, blue!85!black, font=\small] at (2.1, 4.6) {$f(x) = a^x \quad (a > 1)$};
			
			% Curva 0 < a < 1: g(x) = (1/2)^x = 2^(-x)
			\draw[domain=-2.3:3.2, smooth, line width=1.3pt, red!85!black] plot ({\x}, {2^(-\x)});
			\node[left, red!85!black, font=\small] at (-2.1, 4.6) {$g(x) = a^x \quad (0 < a < 1)$};
			
			% Ponto notável comum (0, 1)
			\filldraw[fill=black] (0,1) circle (2.2pt) node[right=3pt, font=\footnotesize] {$(0, 1)$};
			
			% Pontos de referência para x = 1 e x = -1
			\draw[dotted, line width=0.7pt] (1,0) -- (1,2) -- (0,2);
			\filldraw[fill=blue!85!black] (1,2) circle (2pt) node[right=2pt, font=\scriptsize] {$(1, a)$};
			\node[below, font=\footnotesize] at (1,0) {$1$};
			\node[left, font=\footnotesize] at (0,2) {$a$};
			
			\draw[dotted, line width=0.7pt] (-1,0) -- (-1,2) -- (0,2);
			\filldraw[fill=red!85!black] (-1,2) circle (2pt) node[left=2pt, font=\scriptsize] {$(-1, \frac{1}{a})$};
			\node[below, font=\footnotesize] at (-1,0) {$-1$};
		\end{tikzpicture}
	\end{center}
	
	\section{A Assíntota Horizontal}
	\label{subsec:assintota}
	
	Um traço singular da curva exponencial é a existência de uma reta limitante que atrai assintoticamente o gráfico nos extremos do infinito:
	
	\begin{destaque}[Assíntota Horizontal]
		\label{dest:assintota}
		Para a função elementar $f(x) = a^x$, o eixo das abscissas (a reta de equação $y = 0$) é a sua \textbf{assíntota horizontal}.
		
		Se $a>1$, quando $x$ tende ao infinito negativo, $x \to -\infty$, a função se aproxima de zero, mas nunca atinge exatamente 0.
		
		Se $0<a<1$, quando $x$ tende ao infinito , $x \to \infty$, a função também se aproxima de zero, mas, novamente, nunca atinge exatamente 0.
		
		Aqui aperece um conceito matemático que será explorado em breve, no Cálculo I: \textbf{limite}. 
		
		Dizemos que o limite de $a^x$ (para $a>1$), quando $x$ tende ao infinito negativo é zero.
		
		Ou que o limite de $a^x$ (para $0<a<1$), quando $x$ tende ao infinito positivo é zero.
		
		Em uma linguagem matemática formal:
		
		$$\lim_{x \to -\infty} a^x = 0 \quad (\text{se } a > 1) \qquad \text{ou} \qquad \lim_{x \to +\infty} a^x = 0 \quad (\text{se } 0 < a < 1)$$
		
		Na forma transladada $f(x) = k \cdot a^{c(x-x_0)} + d$, a reta assintótica desloca-se para $y = d$.
	\end{destaque}
	
	\section{Interseção com os Eixos Coordenados}
	\label{subsec:intersecoes}
	
	\begin{itemize}
		\item \textbf{Interseção com o eixo das ordenadas ($y$):} Fazendo $x = 0$ na função elementar, obtém-se:
		$$f(0) = a^0 = 1$$
		Logo, qualquer que seja a base $a$, a curva elementar cruza obrigatoriamente o eixo vertical no ponto notável $(0, 1)$. Na forma geral $f(x) = k \cdot a^{cx} + d$, o corte dá-se em $(0, k + d)$.
		\item \textbf{Interseção com o eixo das abscissas ($x$) ou raízes:} Como $a^x > 0$ para todo $x \in \mathbb{R}$, a equação $a^x = 0$ \textbf{não possui solução real}. A curva aproxima-se assintoticamente da reta $y = 0$, mas jamais intercepta o eixo $x$\footnote{Para a forma transladada $f(x) = k \cdot a^{c(x-x_0)} + d$, o cálculo da raíz exige, quase sempre, o uso de uma ferramenta específica, \textbf{logaritmo}, que será estudada em breve.}.
		
		
	\end{itemize}
	
	\chapter{Comportamento Global, Imagem e Sinais}
	\label{sec:caracteristicas_exp}
	
	\section{Domínio e Conjunto Imagem}
	\label{subsec:dominio_imagem}
	
	A função exponencial padrão\footnote{Vamos focar na função exponencial padrão, do tipo $f(x)=a^x$. A forma transladada impacta diretamente em cada característica, mas iremos discutí-las à medida que forem aparecendo.} está definida sobre todo o corpo dos números reais, mas gera apenas valores estritamente positivos:
	
	\begin{destaque}[Domínio e Imagem]
		\label{dest:dom_im}
		Para $f(x) = a^x$ com $a > 0$ e $a \neq 1$:
		$$\text{Dom}(f) = \mathbb{R}$$
		$$\text{Im}(f) = \mathbb{R}^*_+ = \{y \in \mathbb{R} \mid y > 0\} =  {]0, +\infty[}$$
	\end{destaque}
	
	Quando a função sofre reflexão ou translação vertical na forma $g(x) = k \cdot a^x + d$:
	\begin{itemize}
		\item Se $k > 0 \implies \text{Im}(g) = {]d, +\infty[} = \{y \in \mathbb{R} \mid y > d\}$;
		\item Se $k < 0 \implies \text{Im}(g) = {]-\infty, d[} = \{y \in \mathbb{R} \mid y < d\}$.
	\end{itemize}
	
	\section{Crescimento, Decrescimento e Monotonicidade Estrita}
	\label{subsec:monotonia}
	
	Ao contrário da função quadrática, que altera seu sentido de crescimento no vértice, a função exponencial preserva a mesma taxa relativa de variação em todo o domínio:
	
	\begin{destaque}[Monotonicidade]
		\label{dest:monotonia}
		Sejam $x_1, x_2 \in \mathbb{R}$ tais que $x_1 < x_2$:
		\begin{itemize}
			\item Se $a > 1 \implies a^{x_1} < a^{x_2} \quad$ (\textbf{Estritamente Crescente});
			\item Se $0 < a < 1 \implies a^{x_1} > a^{x_2} \quad$ (\textbf{Estritamente Decrescente}).
		\end{itemize}
	\end{destaque}
	
	Uma consequência analítica direta da monotonicidade estrita é a \textbf{injetividade}:
	$$a^{x_1} = a^{x_2} \iff x_1 = x_2$$
	Essa propriedade é a base para a resolução de equações exponenciais elementares e assegura a existência de sua função inversa: a função logarítmica (Em breve!).
	
	\section{Estudo de Sinais}
	\label{subsec:estudo_sinais}
	
	Para a função pura $f(x) = a^x$, o estudo de sinais é imediato em virtude da positividade estrita da potência de base positiva:
	
	\begin{destaque}[Estudo de Sinais para $f(x) = a^x$]
		\label{dest:sinais_pura}
		$$f(x) > 0, \quad \forall x \in \mathbb{R}$$
		A função jamais assume valores nulos ou negativos.
	\end{destaque}
	
	
	\chapter{Uma Sistematização Geral das Funções Exponenciais}
	

	
	\label{sec:sistematizacao_exp}
	
	Apresentamos a síntese de todas as definições e propriedades geométricas desenvolvidas ao longo do texto.
	
	\section{Condição de Existência da Base}
	\begin{destaque}[Condição de Existência]
		$$a \in \mathbb{R}, \quad a > 0 \quad \text{e} \quad a \neq 1$$
	\end{destaque}
	
	\section{Formas de Representação Algébrica}
	\begin{destaque}[Forma Elementar]
		$$f(x) = a^x$$
	\end{destaque}
	
	\begin{destaque}[Forma Geral Parametrizada]
		$$f(x) = k \cdot a^{c(x - x_0)} + d, \quad \text{com } k \neq 0 \text{ e } c \neq 0$$
	\end{destaque}
	
	\section{Domínio e Conjunto Imagem}
	\begin{destaque}[Domínio e Imagem da Forma Elementar]
		$$\text{Dom}(f) = \mathbb{R}$$
		$$\text{Im}(f) = \mathbb{R}^*_+ = {]0, +\infty[}$$
	\end{destaque}
	
	\section{Assíntota Horizontal e Raízes}
	\begin{destaque}[Comportamento Assintótico e Interseção com $x$]
		Para a forma elementar $f(x) = a^x$:
		\begin{itemize}
			\item \textbf{Assíntota Horizontal:} Reta $y = 0$ (o próprio eixo $x$).
			\item \textbf{Raízes Reais:} Inexistentes ($a^x \neq 0, \, \forall x \in \mathbb{R}$).
		\end{itemize}
		Para a forma transladada $f(x) = k \cdot a^{c(x-x_0)} + d$:
		\begin{itemize}
			\item \textbf{Assíntota Horizontal:} Reta $y = d$.
		\end{itemize}
	\end{destaque}
	
	\section{Interseção com o Eixo das Ordenadas $(y)$}
	\begin{destaque}[Interseção Vertical]
		O gráfico elementar intercepta o eixo $y$ no ponto fixo:
		$$(0, 1)$$
		pois $f(0) = a^0 = 1$.
	\end{destaque}
	
	\section{Crescimento e Decrescimento}
	\begin{destaque}[Crescimento e Decrescimento]
		A função elementar $f(x) = a^x$ é:
		\begin{itemize}
			\item \textbf{Estritamente Crescente} em todo o domínio se $a > 1$;
			\item \textbf{Estritamente Decrescente} em todo o domínio se $0 < a < 1$.
		\end{itemize}
	\end{destaque}
	
	\section{Estudo de Sinais}
	\begin{destaque}[Estudo de Sinais]
		Para a forma elementar $f(x) = a^x$:
		\begin{itemize}
			\item $f(x) > 0$ para todo $x \in \mathbb{R}$.
		\end{itemize}
	\end{destaque}
	
	\chapter{Equações Exponenciais: Métodos Fundamentais de Resolução}
	\label{sec:equacoes_exp}
	
	Uma \textbf{equação exponencial} caracteriza-se pela presença da incógnita no expoente de pelo menos uma potência cuja base satisfaz as condições de existência da função exponencial elementar ($a \in \mathbb{R}$, $a > 0$ e $a \neq 1$). Sentenças da forma $2^x = 64$, $3^{2x-1} = 27$ e $4^x - 2^x = 56$ são exemplos típicos dessa classe de problemas. Para resolver esses problemas algébricos e determinar os valores reais que tornam essas igualdades verdadeiras, apoiamo-nos em dois métodos: o \textbf{método da redução a uma base comum} — fundamentado diretamente na injetividade da função exponencial — e o \textbf{método da mudança de variável}, que reduz equações aparentemente mais intrincadas a formas polinomiais clássicas.
	
	\section{Método da Redução a uma Base Comum}
	\label{subsec:metodo_base_comum}
	
	Este método apoia-se na propriedade de monotonicidade estrita demonstrada no Capítulo~\ref{sec:caracteristicas_exp}. Como toda função exponencial elementar $f(x) = a^x$ é estritamente monótona, ela é necessariamente \textbf{injetiva}. Isso garante o seguinte princípio de equivalência:
	
	\begin{destaque}[Propriedade Fundamental da Injetividade]
		Sejam $a \in \mathbb{R}^*_+ \setminus \{1\}$ e $u(x), v(x)$ expressões algébricas com valores reais. Então:
		\[
		a^{u(x)} = a^{v(x)} \iff u(x) = v(x)
		\]
	\end{destaque}
	
	A estratégia consiste em aplicar as propriedades operatórias da potenciação (fatoração em fatores primos, expoentes fracionários e propriedades de potências de potências) em ambos os membros da igualdade até reduzi-los a uma única potência de mesma base. Uma vez estabelecida a igualdade de bases, eliminam-se as bases e igualam-se diretamente os expoentes.
	
	\subsection*{Exemplo Comentado 1: Aplicação direta das propriedades operatórias}
	
	\textbf{Problema:} Determine o conjunto solução da equação:
	\[
	\left( \sqrt{8} \right)^x = \sqrt[5]{32}
	\]
	
	\textbf{Resolução comentada:}
	\begin{enumerate}[label=\alph*)]
		\item \textit{Identificação da base comum:} Como tanto o número $8$ quanto o $32$ são potências de base $2$ ($8 = 2^3$ e $32 = 2^5$), reduzimos os dois membros à base $2$.
		\item \textit{Conversão dos radicais em expoentes racionais:}
		\[
		\sqrt{8} = \sqrt{2^3} = 2^{\frac{3}{2}} \implies \left( \sqrt{8} \right)^x = \left( 2^{\frac{3}{2}} \right)^x = 2^{\frac{3x}{2}}
		\]
		\[
		\sqrt[5]{32} = \sqrt[5]{2^5} = 2^{\frac{5}{5}} = 2^1
		\]
		\item \textit{Aplicação da injetividade:} Igualando as potências de mesma base:
		\[
		2^{\frac{3x}{2}} = 2^1 \iff \frac{3x}{2} = 1
		\]
		\item \textit{Isolamento da incógnita:}
		\[
		3x = 2 \implies x = \frac{2}{3}
		\]
	\end{enumerate}
	Portanto, o conjunto solução é $S = \left\{ \dfrac{2}{3} \right\}$.
	
	\vspace{0.3cm}
	
	\subsection*{Exemplo Comentado 2: Fator comum em evidência}
	
	\textbf{Problema:} Determine o conjunto solução da equação:
	\[
	3^{x+2} + 3^{x-1} - 3^x = 75
	\]
	
	\textbf{Resolução comentada:}
	Quando ocorrem adições ou subtrações de potências com a mesma base e expoentes lineares na incógnita, desmembram-se as potências pelas propriedades de produto e quociente ($a^{x+n} = a^x \cdot a^n$) para colocar a potência elementar em evidência:
	\begin{enumerate}[label=\alph*)]
		\item \textit{Decomposição das potências e fatoração:}
		\[
		3^x \cdot 3^2 + 3^x \cdot 3^{-1} - 3^x \cdot 1 = 75
		\]
		\[
		3^x \left( 9 + \frac{1}{3} - 1 \right) = 75
		\]
		\item \textit{Redução aritmética dos coeficientes:}
		\[
		8 + \frac{1}{3} = \frac{25}{3} \implies 3^x \cdot \frac{25}{3} = 75
		\]
		\item \textit{Isolamento da potência e redução à base comum:}
		\[
		3^x \cdot \frac{25}{3} = 75 \implies 3^x = \frac{75 \cdot 3}{25} = 3 \cdot 3 = 9
		\]
		\[
		3^x = 3^2 \iff x = 2
		\]
	\end{enumerate}
	Portanto, o conjunto solução é $S = \{2\}$.
	
	\section{Método da Mudança de Variável (Incógnita Auxiliar)}
	\label{subsec:metodo_mudanca_variavel}
	
	Em equações onde figuram termos polinomiais com potências da forma $a^{2x}$ associadas a $a^x$ e constantes aditivas, a redução direta de ambos os membros a uma base única não se aplica de modo imediato. Nesses cenários, introduz-se uma variável auxiliar para transformar o problema transcendental em um problema puramente algébrico.
	
	\begin{destaque}[Procedimento de Substituição]
		Dada uma equação na forma:
		\[
		A \cdot a^{2x} + B \cdot a^x + C = 0 \quad (\text{com } A \neq 0)
		\]
		adota-se a substituição $y = a^x$. Decorre que $a^{2x} = (a^x)^2 = y^2$, o que conduz à equação quadrática:
		\[
		A y^2 + B y + C = 0
		\]
		\textbf{Atenção à condição de existência da variável auxiliar:} Como a imagem da função exponencial é restrita aos reais positivos, impõe-se imperativamente:
		\[
		y > 0
		\]
		Portanto, quaisquer raízes não positivas ($y \leqslant 0$) obtidas para $y$ devem ser obrigatoriamente descartadas no retorno à incógnita original $x$.
	\end{destaque}
	
	\subsection*{Exemplo Comentado 3: Resolução via equação de $2^\circ$ grau}
	
	\textbf{Problema:} Determine as raízes reais da equação:
	\[
	9^x - 4 \cdot 3^x + 3 = 0
	\]
	
	\textbf{Resolução comentada:}
	\begin{enumerate}[label=\alph*)]
		\item \textit{Preparação da base de referência:} Escrevemos a base $9$ como potência de base $3$:
		\[
		9^x = (3^2)^x = (3^x)^2
		\]
		Reescrevendo a sentença:
		\[
		(3^x)^2 - 4(3^x) + 3 = 0
		\]
		\item \textit{Introdução da incógnita auxiliar:} Definindo $y = 3^x$ com a condição $y > 0$:
		\[
		y^2 - 4y + 3 = 0
		\]
		\item \textit{Cálculo das raízes quadráticas:}
		Fatorando o trinômio do segundo grau:
		\[
		(y - 1)(y - 3) = 0 \implies y_1 = 1 \quad \text{e} \quad y_2 = 3
		\]
		Como $y_1 > 0$ e $y_2 > 0$, ambas as raízes são admissíveis para a variável auxiliar.
		\item \textit{Retorno à variável original $x$:}
		\begin{itemize}
			\item Para $y = 1$: $3^x = 1 \implies 3^x = 3^0 \iff x = 0$
			\item Para $y = 3$: $3^x = 3 \implies 3^x = 3^1 \iff x = 1$
		\end{itemize}
	\end{enumerate}
	Portanto, o conjunto solução é $S = \{0, 1\}$.
	
	\vspace{0.3cm}
	
	\subsection*{Exemplo Comentado 4: Equações homogêneas resolúveis por divisão de bases}
	
	\textbf{Problema:} Determine a solução real de:
	\[
	9^x + 15^x = 2 \cdot 25^x
	\]
	
	\textbf{Resolução comentada:}
	\begin{enumerate}[label=\alph*)]
		\item \textit{Reconhecimento das potências:}
		Expressando as bases por meio de seus fatores primos ($3$ e $5$):
		\[
		9^x = (3^2)^x = (3^x)^2, \quad 15^x = (3 \cdot 5)^x = 3^x \cdot 5^x, \quad 25^x = (5^2)^x = (5^x)^2
		\]
		Trata-se de uma equação homogênea em relação aos fatores $3^x$ e $5^x$.
		\item \textit{Normalização por divisão:}
		Como $25^x > 0$ para qualquer $x \in \mathbb{R}$, dividimos todos os termos por $25^x$:
		\[
		\frac{9^x}{25^x} + \frac{15^x}{25^x} = \frac{2 \cdot 25^x}{25^x} \implies \left(\frac{9}{25}\right)^x + \left(\frac{15}{25}\right)^x = 2
		\]
		Simplificando as razões fracionárias:
		\[
		\left(\left(\frac{3}{5}\right)^2\right)^x + \left(\frac{3}{5}\right)^x = 2 \implies \left(\frac{3}{5}\right)^{2x} + \left(\frac{3}{5}\right)^x - 2 = 0
		\]
		\item \textit{Aplicação da mudança de variável:}
		Tomando $u = \left(\dfrac{3}{5}\right)^x$ (com a restrição $u > 0$):
		\[
		u^2 + u - 2 = 0 \implies (u + 2)(u - 1) = 0
		\]
		Obtemos $u = 1$ ou $u = -2$. Em virtude da condição $u > 0$, descartamos o valor negativo $u = -2$.
		\item \textit{Resolução para a incógnita original:}
		\[
		\left(\frac{3}{5}\right)^x = 1 \implies \left(\frac{3}{5}\right)^x = \left(\frac{3}{5}\right)^0 \iff x = 0
		\]
	\end{enumerate}
	Portanto, o conjunto solução é $S = \{0\}$.
	
	\section{Exercícios propostos}
	\textbf{Resolva as equações exponenciais a seguir:}
	\begin{multicols}{2}
		\begin{enumerate}
			\item $2^x=64$
			\item $\left( \sqrt{3} \right)^x = \sqrt[3]{81}$
			\item $\left( \sqrt[5]{4} \right) ^x = \dfrac{1}{\sqrt{8}}$
			\item $7^{4x+3}=49$
			\item $2^{x^2-x-16}=16$
			\item $27^{x^2+1}=9^{5x}$
			\item $100 \cdot 10^x = \sqrt[x]{1000^5}$
			\item $3^{2x-1} \cdot 9^{3x+4} = 27^{x+1}$
			\item $(2^x)^{x+4} = 32$
			\item $2^{x-1}+2^{x}+2^{x+1}-2^{x+2}+2^{x+3}=120$
			\item $5^{x-2} - 5^x + 5^{x+1} = 505$
			\item $2 \cdot 4^{x+2}-5 \cdot 4^{x+1}-3 \cdot 2^{2x+1} - 4^x = 20$
			\item $4^x-2^x=56$
			\item $4^x+6^x=2 \cdot 9^x$
			\item $4^{x+1}-9 \cdot 2^x + 2 = 0$
			\item $x^{x^2-5x+6}=1$
			\item $x^x = x$
			\item $x^{x^2-2}=1$
			\item $x^{2x^2-7x+4}=x$
			\item $(x^2- x + 1)^{(2x^2-3x-2)} = 1$
		\end{enumerate} 	
	\end{multicols}
	
\section{Resolução dos exercícios propostos}

Apresentamos as resoluções passo a passo das equações exponenciais propostas, organizadas de forma a explicitar as propriedades operatórias de potências, reduções a bases comuns, mudanças de variáveis e o estudo cuidadoso de bases variáveis no conjunto dos números reais.

\subsection*{Equações e Desenvolvimentos}

\begin{enumerate}[label=\textbf{\arabic*.}]
	
	% Item 1
	\item \textbf{Equação:} $2^x = 64$
	
	Fatorando o segundo membro na base $2$:
	\[
	64 = 2^6
	\]
	Substituindo na equação:
	\[
	2^x = 2^6 \implies x = 6
	\]
	Portanto, $S = \{6\}$.
	
	\vspace{0.5cm}
	
	% Item 2
	\item \textbf{Equação:} $\left( \sqrt{3} \right)^x = \sqrt[3]{81}$
	
	Reescrevemos ambos os membros em potências de base $3$, utilizando expoentes fracionários:
	\[
	\sqrt{3} = 3^{\frac{1}{2}} \implies \left( \sqrt{3} \right)^x = \left( 3^{\frac{1}{2}} \right)^x = 3^{\frac{x}{2}}
	\]
	\[
	81 = 3^4 \implies \sqrt[3]{81} = \sqrt[3]{3^4} = 3^{\frac{4}{3}}
	\]
	Igualando os expoentes:
	\[
	3^{\frac{x}{2}} = 3^{\frac{4}{3}} \implies \frac{x}{2} = \frac{4}{3} \implies x = \frac{8}{3}
	\]
	Portanto, $S = \left\{ \dfrac{8}{3} \right\}$.
	
	\vspace{0.5cm}
	
	% Item 3
	\item \textbf{Equação:} $\left( \sqrt[5]{4} \right)^x = \dfrac{1}{\sqrt{8}}$
	
	Reduzimos os termos à base comum $2$:
	\[
	\sqrt[5]{4} = \sqrt[5]{2^2} = 2^{\frac{2}{5}} \implies \left( \sqrt[5]{4} \right)^x = 2^{\frac{2x}{5}}
	\]
	\[
	\frac{1}{\sqrt{8}} = \frac{1}{\sqrt{2^3}} = \frac{1}{2^{\frac{3}{2}}} = 2^{-\frac{3}{2}}
	\]
	Igualando os expoentes:
	\[
	2^{\frac{2x}{5}} = 2^{-\frac{3}{2}} \implies \frac{2x}{5} = -\frac{3}{2} \implies 4x = -15 \implies x = -\frac{15}{4}
	\]
	Portanto, $S = \left\{ -\dfrac{15}{4} \right\}$.
	
	\vspace{0.5cm}
	
	% Item 4
	\item \textbf{Equação:} $7^{4x+3} = 49$
	
	Fatorando o número $49$ na base $7$:
	\[
	49 = 7^2
	\]
	Substituindo e igualando os expoentes:
	\[
	7^{4x+3} = 7^2 \implies 4x + 3 = 2 \implies 4x = -1 \implies x = -\frac{1}{4}
	\]
	Portanto, $S = \left\{ -\dfrac{1}{4} \right\}$.
	
	\vspace{0.5cm}
	
	% Item 5
	\item \textbf{Equação:} $2^{x^2-x-16} = 16$
	
	Como $16 = 2^4$, temos:
	\[
	2^{x^2-x-16} = 2^4 \implies x^2 - x - 16 = 4 \implies x^2 - x - 20 = 0
	\]
	Resolvendo a equação quadrática:
	\[
	\Delta = (-1)^2 - 4(1)(-20) = 1 + 80 = 81
	\]
	\[
	x = \frac{-(-1) \pm \sqrt{81}}{2(1)} = \frac{1 \pm 9}{2}
	\]
	\[
	x_1 = \frac{1 + 9}{2} = 5 \quad \text{e} \quad x_2 = \frac{1 - 9}{2} = -4
	\]
	Portanto, $S = \{-4, 5\}$.
	
	\vspace{0.5cm}
	
	% Item 6
	\item \textbf{Equação:} $27^{x^2+1} = 9^{5x}$
	
	Expressamos ambos os lados na base $3$:
	\[
	27 = 3^3 \implies \left( 3^3 \right)^{x^2+1} = 3^{3(x^2+1)} = 3^{3x^2+3}
	\]
	\[
	9 = 3^2 \implies \left( 3^2 \right)^{5x} = 3^{2(5x)} = 3^{10x}
	\]
	Igualando os expoentes:
	\[
	3x^2 + 3 = 10x \implies 3x^2 - 10x + 3 = 0
	\]
	Calculando o discriminante:
	\[
	\Delta = (-10)^2 - 4(3)(3) = 100 - 36 = 64
	\]
	\[
	x = \frac{10 \pm \sqrt{64}}{2(3)} = \frac{10 \pm 8}{6}
	\]
	\[
	x_1 = \frac{18}{6} = 3 \quad \text{e} \quad x_2 = \frac{2}{6} = \frac{1}{3}
	\]
	Portanto, $S = \left\{ \dfrac{1}{3}, 3 \right\}$.
	
	\vspace{0.5cm}
	
	% Item 7
	\item \textbf{Equação:} $100 \cdot 10^x = \sqrt[x]{1000^5}$ \quad ($x \neq 0$)
	
	Convertendo os termos para potências de base $10$:
	\[
	100 \cdot 10^x = 10^2 \cdot 10^x = 10^{x+2}
	\]
	\[
	1000^5 = (10^3)^5 = 10^{15} \implies \sqrt[x]{1000^5} = 10^{\frac{15}{x}}
	\]
	Igualando os expoentes:
	\[
	x + 2 = \frac{15}{x}
	\]
	Multiplicando ambos os membros por $x$ (com a condição $x \neq 0$):
	\[
	x(x + 2) = 15 \implies x^2 + 2x - 15 = 0
	\]
	Fatorando o trinômio (Forma Fatorada ou Soma e Produto):
	\[
	(x + 5)(x - 3) = 0 \implies x = -5 \quad \text{ou} \quad x = 3
	\]
	Como ambos os valores são não nulos, são soluções válidas. Logo, $S = \{-5, 3\}$.
	
	\vspace{0.5cm}
	
	% Item 8
	\item \textbf{Equação:} $3^{2x-1} \cdot 9^{3x+4} = 27^{x+1}$
	
	Convertendo para a base $3$:
	\[
	9^{3x+4} = (3^2)^{3x+4} = 3^{6x+8}
	\]
	\[
	27^{x+1} = (3^3)^{x+1} = 3^{3x+3}
	\]
	Multiplicando as potências do primeiro membro:
	\[
	3^{2x-1} \cdot 3^{6x+8} = 3^{(2x-1) + (6x+8)} = 3^{8x+7}
	\]
	Igualando os expoentes:
	\[
	8x + 7 = 3x + 3 \implies 5x = -4 \implies x = -\frac{4}{5}
	\]
	Portanto, $S = \left\{ -\dfrac{4}{5} \right\}$.
	
	\vspace{0.5cm}
	
	% Item 9
	\item \textbf{Equação:} $(2^x)^{x+4} = 32$
	
	Aplicando a propriedade de potência de potência e fatorando $32$:
	\[
	2^{x(x+4)} = 2^5 \implies 2^{x^2+4x} = 2^5
	\]
	Igualando os expoentes:
	\[
	x^2 + 4x = 5 \implies x^2 + 4x - 5 = 0
	\]
	Fatorando o trinômio:
	\[
	(x + 5)(x - 1) = 0 \implies x = -5 \quad \text{ou} \quad x = 1
	\]
	Portanto, $S = \{-5, 1\}$.
	
	\vspace{0.5cm}
	
	% Item 10
	\item \textbf{Equação:} $2^{x-1} + 2^x + 2^{x+1} - 2^{x+2} + 2^{x+3} = 120$
	
	Colocamos o termo comum $2^x$ em evidência:
	\[
	2^x \left( 2^{-1} + 1 + 2^1 - 2^2 + 2^3 \right) = 120
	\]
	Calculando a soma dos coeficientes numéricos:
	\[
	\frac{1}{2} + 1 + 2 - 4 + 8 = \frac{1}{2} + 7 = \frac{15}{2}
	\]
	Substituindo de volta na equação:
	\[
	2^x \cdot \frac{15}{2} = 120 \implies 2^x = \frac{120 \cdot 2}{15} \implies 2^x = 8 \cdot 2 \implies 2^x = 16
	\]
	Como $16 = 2^4$, temos $x = 4$. Portanto, $S = \{4\}$.
	
	\vspace{0.5cm}
	
	% Item 11
	\item \textbf{Equação:} $5^{x-2} - 5^x + 5^{x+1} = 505$
	
	Evidenciamos o fator comum $5^x$:
	\[
	5^x \left( 5^{-2} - 1 + 5^1 \right) = 505
	\]
	Calculando o termo entre parênteses:
	\[
	\frac{1}{25} - 1 + 5 = 4 + \frac{1}{25} = \frac{101}{25}
	\]
	Isolando $5^x$:
	\[
	5^x \cdot \frac{101}{25} = 505 \implies 5^x = \frac{505 \cdot 25}{101}
	\]
	Observando que $505 / 101 = 5$:
	\[
	5^x = 5 \cdot 25 \implies 5^x = 125 \implies 5^x = 5^3 \implies x = 3
	\]
	Portanto, $S = \{3\}$.
	
	\vspace{0.5cm}
	
	% Item 12
	\item \textbf{Equação:} $2 \cdot 4^{x+2} - 5 \cdot 4^{x+1} - 3 \cdot 2^{2x+1} - 4^x = 20$
	
	Expressamos todos os termos em função de $4^x$, lembrando que $2^{2x+1} = 2^{2x} \cdot 2^1 = (2^2)^x \cdot 2 = 2 \cdot 4^x$:
	\[
	2 \cdot 4^{x+2} = 2 \cdot 4^2 \cdot 4^x = 32 \cdot 4^x
	\]
	\[
	5 \cdot 4^{x+1} = 5 \cdot 4^1 \cdot 4^x = 20 \cdot 4^x
	\]
	\[
	3 \cdot 2^{2x+1} = 3 \cdot (2 \cdot 4^x) = 6 \cdot 4^x
	\]
	Substituindo na expressão original:
	\[
	32 \cdot 4^x - 20 \cdot 4^x - 6 \cdot 4^x - 1 \cdot 4^x = 20
	\]
	Reduzindo os termos semelhantes:
	\[
	(32 - 20 - 6 - 1) \cdot 4^x = 20 \implies 5 \cdot 4^x = 20 \implies 4^x = 4 \implies x = 1
	\]
	Portanto, $S = \{1\}$.
	
	\vspace{0.5cm}
	
	% Item 13
	\item \textbf{Equação:} $4^x - 2^x = 56$
	
	Reescrevemos $4^x$ como $(2^2)^x = (2^x)^2$:
	\[
	(2^x)^2 - 2^x - 56 = 0
	\]
	Fazemos a mudança de variável $y = 2^x$ (com a condição $y > 0$):
	\[
	y^2 - y - 56 = 0
	\]
	Fatorando o trinômio:
	\[
	(y - 8)(y + 7) = 0 \implies y = 8 \quad \text{ou} \quad y = -7
	\]
	Como $y = 2^x > 0$, descartamos a raiz negativa $y = -7$. Para a raiz positiva:
	\[
	2^x = 8 \implies 2^x = 2^3 \implies x = 3
	\]
	Portanto, $S = \{3\}$.
	
	\vspace{0.5cm}
	
		% Item 14
	\item \textbf{Equação:} $4^x + 6^x = 2 \cdot 9^x$
	
	Notamos que as bases envolvidas podem ser expressas em função dos fatores primos $2$ e $3$:
	\[
	4^x = (2^2)^x = (2^x)^2, \quad 6^x = (2 \cdot 3)^x = 2^x \cdot 3^x, \quad 9^x = (3^2)^x = (3^x)^2
	\]
	Como $9^x > 0$ para todo $x \in \mathbb{R}$, dividimos a equação inteira por $9^x$:
	\[
	\frac{4^x}{9^x} + \frac{6^x}{9^x} = \frac{2 \cdot 9^x}{9^x}
	\]
	Simplificando as razões:
	\[
	\left( \frac{4}{9} \right)^x + \left( \frac{6}{9} \right)^x = 2 \implies \left( \left( \frac{2}{3} \right)^2 \right)^x + \left( \frac{2}{3} \right)^x = 2
	\]
	\[
	\left( \frac{2}{3} \right)^{2x} + \left( \frac{2}{3} \right)^x - 2 = 0
	\]
	Introduzimos a variável auxiliar $u = \left( \dfrac{2}{3} \right)^x$, com $u > 0$:
	\[
	u^2 + u - 2 = 0
	\]
	Fatorando o trinômio:
	\[
	(u + 2)(u - 1) = 0 \implies u = 1 \quad \text{ou} \quad u = -2
	\]
	Como $u > 0$, descartamos $u = -2$. Retornando à variável original para $u = 1$:
	\[
	\left( \frac{2}{3} \right)^x = 1 \implies \left( \frac{2}{3} \right)^x = \left( \frac{2}{3} \right)^0 \implies x = 0
	\]
	Portanto, $S = \{0\}$.
	
	\vspace{0.5cm}
	
	% Item 15
	\item \textbf{Equação:} $4^{x+1} - 9 \cdot 2^x + 2 = 0$
	
	Transformamos a base $4$ na base $2$:
	\[
	4^{x+1} = 4^1 \cdot (2^2)^x = 4 \cdot (2^x)^2
	\]
	Substituindo na equação:
	\[
	4(2^x)^2 - 9(2^x) + 2 = 0
	\]
	Adotando a variável auxiliar $y = 2^x$ ($y > 0$):
	\[
	4y^2 - 9y + 2 = 0
	\]
	Calculando o discriminante:
	\[
	\Delta = (-9)^2 - 4(4)(2) = 81 - 32 = 49
	\]
	\[
	y = \frac{9 \pm \sqrt{49}}{2(4)} = \frac{9 \pm 7}{8} \implies y_1 = \frac{16}{8} = 2 \quad \text{e} \quad y_2 = \frac{2}{8} = \frac{1}{4}
	\]
	Retornando à variável $x$:
	\begin{itemize}
		\item Para $y = 2$: $2^x = 2^1 \implies x = 1$
		\item Para $y = \frac{1}{4}$: $2^x = 2^{-2} \implies x = -2$
	\end{itemize}
	Portanto, $S = \{-2, 1\}$.
	
	\vspace{0.5cm}
	
	% Item 16
	\item \textbf{Equação:} $x^{x^2-5x+6} = 1$
	
	A igualdade $a^b = 1$ no conjunto dos números reais ocorre em três casos:
	\begin{enumerate}[label=\alph*)]
		\item \textbf{Expoente nulo com base não nula:}
		\[
		x^2 - 5x + 6 = 0 \quad \text{e} \quad x \neq 0
		\]
		\[
		(x - 2)(x - 3) = 0 \implies x = 2 \quad \text{ou} \quad x = 3
		\]
		Ambos os valores satisfazem $x \neq 0$.
		
		\item \textbf{Base igual a $1$:}
		\[
		x = 1
		\]
		Verificação direta: $1^{1^2 - 5(1) + 6} = 1^2 = 1$ (verdadeiro).
		
		\item \textbf{Base igual a $-1$ com expoente inteiro par:}
		\[
		x = -1
		\]
		Avaliando o expoente em $x = -1$:
		\[
		(-1)^2 - 5(-1) + 6 = 1 + 5 + 6 = 12
		\]
		Como $12$ é um número inteiro par, $(-1)^{12} = 1$ (verdadeiro).
	\end{enumerate}
	Reunindo os três casos, temos $S = \{-1, 1, 2, 3\}$.
	
	\vspace{0.5cm}
	
	% Item 17
	\item \textbf{Equação:} $x^x = x$
	
	Reescrevemos a equação na forma $x^x = x^1$. Analisamos os casos no conjunto dos reais:
	\begin{enumerate}[label=\alph*)]
		\item \textbf{Expoentes iguais com base diferente de $0$, $1$ e $-1$:}
		\[
		x = 1
		\]
		Contudo, a base $x = 1$ constitui um caso próprio examinado a seguir.
		
		\item \textbf{Base igual a $1$:}
		\[
		x = 1 \implies 1^1 = 1
		\]
		Identidade confirmada; logo, $x = 1$ é solução.
		
		\item \textbf{Base igual a $0$ com expoente estritamente positivo:}
		\[
		x = 0
		\]
		Substituindo na expressão, teríamos $0^0$, que é uma indeterminação nos reais. Logo, $x = 0$ não é solução.
		
		\item \textbf{Base igual a $-1$ com expoente inteiro ímpar:}
		\[
		x = -1
		\]
		O expoente é $x = -1$, que é inteiro e ímpar:
		\[
		(-1)^{-1} = \frac{1}{(-1)^1} = -1
		\]
		A igualdade é satisfeita, confirmando que $x = -1$ é solução.
	\end{enumerate}
	Reunindo os valores válidos, obtemos $S = \{-1, 1\}$.
	
	\vspace{0.5cm}
	
	% Item 18
	\item \textbf{Equação:} $x^{x^2-2} = 1$
	
	Analisamos as condições para que uma potência resulte em $1$ nos reais:
	\begin{enumerate}[label=\alph*)]
		\item \textbf{Expoente nulo com base não nula:}
		\[
		x^2 - 2 = 0 \quad \text{e} \quad x \neq 0
		\]
		\[
		x^2 = 2 \implies x = \sqrt{2} \quad \text{ou} \quad x = -\sqrt{2}
		\]
		Como $\pm\sqrt{2} \neq 0$, ambos são soluções.
		
		\item \textbf{Base igual a $1$:}
		\[
		x = 1 \implies 1^{1^2 - 2} = 1^{-1} = 1
		\]
		Identidade confirmada; logo, $x = 1$ é solução.
		
		\item \textbf{Base igual a $-1$ com expoente inteiro par:}
		\[
		x = -1
		\]
		Avaliando o expoente em $x = -1$:
		\[
		(-1)^2 - 2 = 1 - 2 = -1
		\]
		Como $-1$ é ímpar, $(-1)^{-1} = -1 \neq 1$. Portanto, $x = -1$ não é solução.
	\end{enumerate}
	Reunindo as soluções encontradas, temos $S = \{-\sqrt{2}, 1, \sqrt{2}\}$.
	
	\vspace{0.5cm}
	
	% Item 19
	\item \textbf{Equação:} $x^{2x^2-7x+4} = x$
	
	Reescrevemos o segundo membro como $x^1$. Analisamos os casos admissíveis nos reais:
	\begin{enumerate}[label=\alph*)]
		\item \textbf{Expoentes iguais com base diferente de $0$, $1$ e $-1$:}
		\[
		2x^2 - 7x + 4 = 1 \implies 2x^2 - 7x + 3 = 0
		\]
		\[
		\Delta = (-7)^2 - 4(2)(3) = 49 - 24 = 25
		\]
		\[
		x = \frac{7 \pm 5}{4} \implies x = 3 \quad \text{ou} \quad x = \frac{1}{2}
		\]
		Ambos os valores diferem de $0$, $1$ e $-1$, sendo soluções válidas.
		
		\item \textbf{Base igual a $1$:}
		\[
		x = 1 \implies 1^{2(1)^2 - 7(1) + 4} = 1^{-1} = 1
		\]
		A igualdade se mantém válida.
		
		\item \textbf{Base igual a $0$ com expoente estritamente positivo:}
		\[
		x = 0
		\]
		Calculando o expoente: $2(0)^2 - 7(0) + 4 = 4 > 0$. Logo, $0^4 = 0$, confirmando a solução.
		
		\item \textbf{Base igual a $-1$ com expoente inteiro ímpar:}
		\[
		x = -1
		\]
		Calculando o expoente: $2(-1)^2 - 7(-1) + 4 = 2(1) + 7 + 4 = 13$. Como $13$ é inteiro e ímpar, $(-1)^{13} = -1$, satisfazendo a igualdade.
	\end{enumerate}
	Reunindo todas as possibilidades, obtemos $S = \left\{ -1, 0, \dfrac{1}{2}, 1, 3 \right\}$.
	
	\vspace{0.5cm}
	
	% Item 20 (Nova: (x^2-x+1)^{2x^2-3x-2} = 1)
	\item \textbf{Equação:} $(x^2 - x + 1)^{(2x^2 - 3x - 2)} = 1$
	
	Para uma igualdade do tipo $[f(x)]^{g(x)} = 1$, analisamos as três situações nos reais:
	\begin{enumerate}[label=\alph*)]
		\item \textbf{Expoente nulo com base não nula:}
		\[
		2x^2 - 3x - 2 = 0 \quad \text{e} \quad x^2 - x + 1 \neq 0
		\]
		Calculando as raízes do expoente:
		\[
		\Delta = (-3)^2 - 4(2)(-2) = 9 + 16 = 25
		\]
		\[
		x = \frac{3 \pm 5}{4} \implies x = 2 \quad \text{ou} \quad x = -\frac{1}{2}
		\]
		Verificamos a base $f(x) = x^2 - x + 1$:
		\begin{itemize}
			\item Para $x = 2$: $2^2 - 2 + 1 = 3 \neq 0$.
			\item Para $x = -\frac{1}{2}$: $\left(-\frac{1}{2}\right)^2 - \left(-\frac{1}{2}\right) + 1 = \frac{1}{4} + \frac{1}{2} + 1 = \frac{7}{4} \neq 0$.
		\end{itemize}
		Como a base não se anula, $x = 2$ e $x = -\dfrac{1}{2}$ são soluções.
		
		\item \textbf{Base igual a $1$:}
		\[
		x^2 - x + 1 = 1 \implies x^2 - x = 0 \implies x(x - 1) = 0 \implies x = 0 \quad \text{ou} \quad x = 1
		\]
		Para a base igual a $1$, a potência $1^{g(x)} = 1$ é satisfeita para qualquer expoente real:
		\begin{itemize}
			\item Para $x = 0$: $g(0) = -2$, e $1^{-2} = 1$.
			\item Para $x = 1$: $g(1) = 2(1)^2 - 3(1) - 2 = -3$, e $1^{-3} = 1$.
		\end{itemize}
		Logo, $x = 0$ e $x = 1$ são soluções.
		
		\item \textbf{Base igual a $-1$ com expoente inteiro par:}
		\[
		x^2 - x + 1 = -1 \implies x^2 - x + 2 = 0
		\]
		Calculando o discriminante da base:
		\[
		\Delta = (-1)^2 - 4(1)(2) = 1 - 8 = -7 < 0
		\]
		Como a equação não admite raízes reais, a base nunca assume o valor $-1$ nos reais.
	\end{enumerate}
	Reunindo todas as soluções reais, obtemos $S = \left\{ -\dfrac{1}{2}, 0, 1, 2 \right\}$.
	
\end{enumerate}

	\chapter{Inequações Exponenciais}
	\label{sec:inequacoes_exp}
	
	Uma \textbf{inequação exponencial} é uma desigualdade condicionada caracterizada pela presença da incógnita no expoente de pelo menos uma potência cuja base satisfaz as condições de existência da função exponencial elementar ($a \in \mathbb{R}$, $a > 0$ e $a \neq 1$), expressa por meio das relações de ordem $<$, $\leqslant$, $>$ ou $\geqslant$. Enquanto na resolução de equações a propriedade central utilizada é a injetividade ($a^u = a^v \iff u = v$), nas inequações o elemento analítico determinante é a \textbf{monotonicidade estrita} da base. Assim, ao transpor a desigualdade das potências para os seus respectivos expoentes, o sentido do sinal de desigualdade é preservado quando a base for estritamente maior que $1$ e invertido quando a base estiver compreendida entre $0$ e $1$.
	
	\section{Fundamento Teórico e Propriedades de Monotonia}
	\label{subsec:propriedade_inequacoes}
	
	A validade da passagem de uma desigualdade entre potências de mesma base para uma desigualdade entre seus expoentes decorre diretamente do estudo do crescimento e decrescimento da função exponencial $f(x) = a^x$:
	
	\begin{destaque}[Princípio Fundamental das Inequações Exponenciais]
		Sejam $u(x)$ e $v(x)$ funções reais e $a \in \mathbb{R}^*_+ \setminus \{1\}$:
		\begin{enumerate}
			\item \textbf{Se a base for maior do que um ($a > 1$):} A função exponencial é estritamente crescente, logo \textbf{mantém-se} o sentido da desigualdade:
			\[
			a^{u(x)} > a^{v(x)} \iff u(x) > v(x) \qquad \text{e} \qquad a^{u(x)} < a^{v(x)} \iff u(x) < v(x)
			\]
			\item \textbf{Se a base estiver entre zero e um ($0 < a < 1$):} A função exponencial é estritamente decrescente, logo \textbf{inverte-se} o sentido da desigualdade:
			\[
			a^{u(x)} > a^{v(x)} \iff u(x) < v(x) \qquad \text{e} \qquad a^{u(x)} < a^{v(x)} \iff u(x) > v(x)
			\]
		\end{enumerate}
	\end{destaque}
	
	\section{Exemplos Comentados}
	\label{subsec:exemplos_inequacoes}
	
	\subsection*{Exemplo Comentado 1: Base maior do que $1$ ($a > 1$)}
	
	\textbf{Problema:} Resolva a inequação $2^{3x - 1} \geqslant 32$.
	
	\textbf{Resolução comentada:}
	\begin{enumerate}[label=\alph*)]
		\item \textit{Redução à base comum:} $32 = 2^5$, logo:
		\[
		2^{3x - 1} \geqslant 2^5
		\]
		\item \textit{Análise da base:} A base é $a = 2$. Como $2 > 1$, a função é estritamente crescente, preservando-se a relação de ordem nos expoentes:
		\[
		3x - 1 \geqslant 5 \implies 3x \geqslant 6 \implies x \geqslant 2
		\]
	\end{enumerate}
	Portanto, o conjunto solução é $S = \{x \in \mathbb{R} \mid x \geqslant 2\} = [2, +\infty[$.
	
	\vspace{0.3cm}
	
	\subsection*{Exemplo Comentado 2: Base entre $0$ e $1$ ($0 < a < 1$)}
	
	\textbf{Problema:} Resolva a inequação $\left(\dfrac{1}{3}\right)^{2x + 5} < \dfrac{1}{27}$.
	
	\textbf{Resolução comentada:}
	\begin{enumerate}[label=\alph*)]
		\item \textit{Redução à base comum:} $\dfrac{1}{27} = \left(\dfrac{1}{3}\right)^3$, reescrevendo a sentença:
		\[
		\left(\frac{1}{3}\right)^{2x + 5} < \left(\frac{1}{3}\right)^3
		\]
		\item \textit{Análise da base:} A base é $a = \dfrac{1}{3}$. Como $0 < \dfrac{1}{3} < 1$, a função é estritamente decrescente. Ao passar a desigualdade para os expoentes, \textbf{inverte-se} o sinal:
		\[
		2x + 5 > 3 \implies 2x > -2 \implies x > -1
		\]
	\end{enumerate}
	Portanto, $S = \{x \in \mathbb{R} \mid x > -1\} = {]-1, +\infty [}$.
	
	\vspace{0.3cm}
	
	\subsection*{Exemplo Comentado 3: Mudança de variável com estudo de sinais}
	
	\textbf{Problema:} Resolva em $\mathbb{R}$ a inequação $4^x - 5 \cdot 2^x + 4 \leqslant 0$.
	
	\textbf{Resolução comentada:}
	\begin{enumerate}[label=\alph*)]
		\item \textit{Preparação e mudança de variável:} Reescrevemos $4^x = (2^x)^2$:
		
		\[ (2^x)^2-5 \cdot 2^x + 4 \leqslant 0 \]
		
		
		\item Introduzimos a incógnita auxiliar $y = 2^x$ (com a restrição obrigatória $y > 0$):
		\[
		y^2 - 5y + 4 \leqslant 0
		\]
		\item \textit{Resolução da inequação quadrática auxiliar:} As raízes da equação associada são dadas por $(y - 1)(y - 4) = 0$, isto é, $y = 1$ ou $y = 4$. Por ser uma parábola com concavidade voltada para cima, o trinômio é não positivo no intervalo entre as raízes:
		\[
		1 \leqslant y \leqslant 4
		\]
		Ambos os extremos satisfazem a condição de positividade $y > 0$.
		\item \textit{Retorno à variável original:}
		\[
		1 \leqslant 2^x \leqslant 4 \implies 2^0 \leqslant 2^x \leqslant 2^2
		\]
		Como a base $a = 2 > 1$, mantém-se o ordenamento entre os expoentes:
		\[
		0 \leqslant x \leqslant 2
		\]
	\end{enumerate}
	Portanto, $S = \{x \in \mathbb{R} \mid 0 \leqslant x \leqslant 2\} = [0, 2]$.
	
	\section{Exercícios Propostos}
	\label{subsec:propostos_inequacoes}
	
	\textbf{Resolva as seguintes inequações exponenciais em $\mathbb{R}$:}
	\begin{multicols}{2}
		\begin{enumerate}
			\item $3^{x+1} > 81$
			\item $\left( \dfrac{1}{2} \right)^{3x-2} \leqslant \dfrac{1}{16}$
			\item $5^{x^2 - 2x} < 125$
			\item $\left( \dfrac{2}{5} \right)^{x^2 - 4} \geqslant 1$
			\item $2^{x-1} + 2^x + 2^{x+1} \leqslant 56$
			\item $9^x - 10 \cdot 3^x + 9 > 0$
			\item $\left( \dfrac{1}{4} \right)^x - 3 \cdot \left( \dfrac{1}{2} \right)^x - 4 < 0$
			\item $4^x - 3 \cdot 2^{x+1} + 8 \leqslant 0$
		\end{enumerate}
	\end{multicols}
	
	\newpage
	
	\section{Resolução dos Exercícios Propostos}
	\label{subsec:resolucao_inequacoes}
	
	\begin{enumerate}[label=\textbf{\arabic*.}]
		
		% Item 1
		\item \textbf{Inequação:} $3^{x+1} > 81$
		
		Fatorando $81$ na base $3$:
		\[
		81 = 3^4 \implies 3^{x+1} > 3^4
		\]
		Como a base $3 > 1$, a função é estritamente crescente, preservando o sinal da desigualdade:
		\[
		x + 1 > 4 \implies x > 3
		\]
		Portanto, $S = \{x \in \mathbb{R} \mid x > 3\} ={]3, +\infty[}$.
		
		\vspace{0.4cm}
		
		% Item 2
		\item \textbf{Inequação:} $\left( \dfrac{1}{2} \right)^{3x-2} \leqslant \dfrac{1}{16}$
		
		Expressamos o segundo membro na base $\dfrac{1}{2}$:
		\[
		\frac{1}{16} = \left(\frac{1}{2}\right)^4 \implies \left(\frac{1}{2}\right)^{3x-2} \leqslant \left(\frac{1}{2}\right)^4
		\]
		Como a base satisfaz $0 < \dfrac{1}{2} < 1$, a função é estritamente decrescente. Invertendo a desigualdade entre os expoentes:
		\[
		3x - 2 \geqslant 4 \implies 3x \geqslant 6 \implies x \geqslant 2
		\]
		Portanto, $S = \{x \in \mathbb{R} \mid x \geqslant 2\} = [2, +\infty[$.
		
		\vspace{0.4cm}
		
		% Item 3
		\item \textbf{Inequação:} $5^{x^2 - 2x} < 125$
		
		Escrevendo $125$ como potência de base $5$:
		\[
		125 = 5^3 \implies 5^{x^2 - 2x} < 5^3
		\]
		Sendo a base $5 > 1$, mantemos o sentido da relação de ordem:
		\[
		x^2 - 2x < 3 \implies x^2 - 2x - 3 < 0
		\]
		Determinando as raízes do polinômio quadrático associado:
		\[
		(x - 3)(x + 1) = 0 \implies x_1 = -1 \quad \text{e} \quad x_2 = 3
		\]
		Como o coeficiente dominante de $x^2$ é positivo, a parábola assume valores negativos no intervalo interior entre as raízes:
		\[
		-1 < x < 3
		\]
		Portanto, $S = \{x \in \mathbb{R} \mid -1 < x < 3\} ={]-1, 3[}$.
		
		\vspace{0.4cm}
		
		% Item 4
		\item \textbf{Inequação:} $\left( \dfrac{2}{5} \right)^{x^2 - 4} \geqslant 1$
		
		Expressamos o número $1$ na base correspondente: $1 = \left(\dfrac{2}{5}\right)^0$:
		\[
		\left(\frac{2}{5}\right)^{x^2 - 4} \geqslant \left(\frac{2}{5}\right)^0
		\]
		Como $0 < \dfrac{2}{5} < 1$, a função exponencial correspondente é estritamente decrescente, o que impõe a inversão do sinal:
		\[
		x^2 - 4 \leqslant 0 \implies (x - 2)(x + 2) \leqslant 0
		\]
		Analisando o sinal da função quadrática cujas raízes são $\pm 2$:
		\[
		-2 \leqslant x \leqslant 2
		\]
		Portanto, $S = \{x \in \mathbb{R} \mid -2 \leqslant x \leqslant 2\} = [-2, 2]$.
		
		\vspace{0.4cm}
		
		% Item 5
		\item \textbf{Inequação:} $2^{x-1} + 2^x + 2^{x+1} \leqslant 56$
		
		Colocamos $2^x$ em evidência no primeiro membro:
		\[
		2^x \left( 2^{-1} + 1 + 2^1 \right) \leqslant 56 \implies 2^x \left( \frac{1}{2} + 3 \right) \leqslant 56
		\]
		\[
		2^x \cdot \frac{7}{2} \leqslant 56 \implies 2^x \leqslant \frac{56 \cdot 2}{7} \implies 2^x \leqslant 8 \cdot 2 \implies 2^x \leqslant 16
		\]
		Escrevendo $16$ na base $2$:
		\[
		2^x \leqslant 2^4
		\]
		Como a base $2 > 1$, mantém-se o sinal da desigualdade:
		\[
		x \leqslant 4
		\]
		Portanto, $S = \{x \in \mathbb{R} \mid x \leqslant 4\} = {]-\infty, 4]}$.
		
		\vspace{0.4cm}
		
		% Item 6
		\item \textbf{Inequação:} $9^x - 10 \cdot 3^x + 9 > 0$
		
		Reescrevemos a inequação com a base $3$:
		\[
		(3^x)^2 - 10 \cdot 3^x + 9 > 0
		\]
		Adotando a substituição $y = 3^x$ com a condição fundamental $y > 0$:
		\[
		y^2 - 10y + 9 > 0
		\]
		Fatorando a expressão:
		\[
		(y - 1)(y - 9) > 0
		\]
		A inequação é satisfeita para $y < 1$ ou $y > 9$. Considerando a condição $y > 0$, dividimos em dois casos:
		\begin{itemize}
			\item $0 < y < 1 \implies 0 < 3^x < 3^0$. Como $3^x > 0$ para todo $x$, resta $3^x < 3^0 \iff x < 0$.
			\item $y > 9 \implies 3^x > 3^2 \iff x > 2$.
		\end{itemize}
		Portanto, $S = \{x \in \mathbb{R} \mid x < 0 \text{ ou } x > 2\} = {]-\infty, 0[} \cup {]2, +\infty[}$.
		
		\vspace{0.4cm}
		
		% Item 7
		\item \textbf{Inequação:} $\left( \dfrac{1}{4} \right)^x - 3 \cdot \left( \dfrac{1}{2} \right)^x - 4 < 0$
		
		Reconhecendo que $\left(\dfrac{1}{4}\right)^x = \left(\left(\dfrac{1}{2}\right)^2\right)^x = \left[\left(\dfrac{1}{2}\right)^x\right]^2$, definimos $y = \left(\dfrac{1}{2}\right)^x$ com $y > 0$:
		\[
		y^2 - 3y - 4 < 0
		\]
		Fatorando o trinômio quadrático:
		\[
		(y - 4)(y + 1) < 0 \implies -1 < y < 4
		\]
		Como $y = \left(\dfrac{1}{2}\right)^x > 0$ para todo $x \in \mathbb{R}$, a condição $-1 < y$ é identicamente satisfeita. Resta resolver:
		\[
		y < 4 \implies \left(\frac{1}{2}\right)^x < 4
		\]
		Expressando $4$ na base $\dfrac{1}{2}$:
		\[
		4 = 2^2 = \left( \left(\frac{1}{2}\right)^{-1} \right)^2 = \left(\frac{1}{2}\right)^{-2}
		\]
		Substituindo na desigualdade:
		\[
		\left(\frac{1}{2}\right)^x < \left(\frac{1}{2}\right)^{-2}
		\]
		Como a base $0 < \dfrac{1}{2} < 1$, inverte-se o sinal da desigualdade:
		\[
		x > -2
		\]
		Portanto, $S = \{x \in \mathbb{R} \mid x > -2\} = {]-2, +\infty[}$.
		
		\vspace{0.4cm}
		
		% Item 8
		\item \textbf{Inequação:} $4^x - 3 \cdot 2^{x+1} + 8 \leqslant 0$
		
		Desenvolvemos as potências em base $2$:
		\[
		2^{x+1} = 2^1 \cdot 2^x = 2 \cdot 2^x \implies 3 \cdot 2^{x+1} = 6 \cdot 2^x
		\]
		\[
		(2^x)^2 - 6 \cdot 2^x + 8 \leqslant 0
		\]
		Fazendo $y = 2^x$ (com $y > 0$):
		\[
		y^2 - 6y + 8 \leqslant 0
		\]
		Fatorando a expressão:
		\[
		(y - 2)(y - 4) \leqslant 0 \implies 2 \leqslant y \leqslant 4
		\]
		Retornando à variável original $x$:
		\[
		2 \leqslant 2^x \leqslant 4 \implies 2^1 \leqslant 2^x \leqslant 2^2
		\]
		Como a base $2 > 1$, preserva-se o sentido da ordenação:
		\[
		1 \leqslant x \leqslant 2
		\]
		Portanto, $S = \{x \in \mathbb{R} \mid 1 \leqslant x \leqslant 2\} = [1, 2]$.
		
	\end{enumerate}
\end{document}